\documentclass[11pt]{article}
\usepackage[margin=0.85in]{geometry}
\usepackage{setspace}
\usepackage{times}
\usepackage{amsmath,amssymb}
\usepackage[colorlinks,citecolor=blue,linkcolor=blue,urlcolor=blue]{hyperref}
\urlstyle{same}

\begin{document}
\singlespacing

\begin{center}
    {\footnotesize\scshape Manuscript submitted to: \textbf{Terrorism and Political Violence} --- Routledge / Taylor \& Francis\par}
    {\scriptsize Editor-in-Chief: Prof. James J.F. Forest \quad $\cdot$ \quad Review Mode: Double-Anonymous Peer Review\par}
    \vspace{0.35cm}
    
    {\Large\bfseries Identifying Rational Types in Unknown Environments under an Indirect Dynamic Mechanism\par}
    
    \vspace{0.35cm}
    
    {\large\textbf{Prof. Humberto Bernal, Ph.D.}\par}
    \vspace{0.12cm}
    {\small Economics Program, Universidad Colegio Mayor de Cundinamarca;\\
    Ph.D. in Economics, Universidad de los Andes, Bogot\'a, Colombia\par}
    \vspace{0.15cm}
    {\footnotesize Corresponding Author $\cdot$ E-mails: \href{mailto:hbernal@uniandes.edu.co}{hbernal@uniandes.edu.co}, \href{mailto:hbernalc@universidadmayor.edu.co}{hbernalc@universidadmayor.edu.co} $\cdot$ ORCID: \href{https://orcid.org/0000-0002-4864-1726}{0000-0002-4864-1726}\par}
    \vspace{0.12cm}
    {\footnotesize\textbf{Submission Date:} October 2026\par}
\end{center}

\vspace{0.2cm}

\begin{abstract}
\noindent This paper develops an indirect dynamic mechanism to identify rational types in unknown environments characterized by observational pooling and asymmetric information. The central theoretical contribution establishes the Rational Type Identification Theorem, proving that the principal's Bayesian posterior beliefs converge almost surely to the agent's true latent type along prolonged interactions without direct revelation reports. This asymptotic identification is operationalized by the Mano de Dios--Guadalupe (MDG) process, an endogenous maximum-entropy implementation noise kernel that ensures full support across the decision space within a competing-risks architecture. A complementary Implementability Theorem guarantees sequential execution in Perfect Bayesian Equilibrium. Microeconometric calibration and deep-learning computational simulations of kidnapping for ransom in Colombia validate rapid type separation and policy feasibility across Colombian insurgent groups and criminal syndicates, including the Revolutionary Armed Forces of Colombia (FARC), the National Liberation Army (ELN), paramilitaries (PAR), and common delinquency (DC).
\end{abstract}

\vspace{0.25cm}

\noindent \textbf{Keywords:} terrorism; political violence; indirect mechanism design; dynamic screening; Bayesian learning; kidnapping for ransom; insurgent groups\\[0.25em]
\noindent \textbf{JEL Classifications:} C73, D82, D86, K42

\vspace{0.3cm}
\noindent\rule{\textwidth}{0.4pt}
\vspace{0.15cm}

\noindent {\large\textbf{Statements and Declarations:}}
\begin{itemize}\footnotesize
    \setlength{\itemsep}{0.2em}
    \item \textbf{Funding:} The author declares that no funds, grants, or other financial support were received for the preparation of this manuscript.
    \item \textbf{Conflict of Interest / Competing Interests:} The author declares that he has no conflict of interest. The author has no relevant financial or non-financial competing interests to disclose.
    \item \textbf{Data and Code Availability Statement:} The datasets generated and analysed during the current study, as well as the replication scripts, empirical estimations, neural network architectures, and interactive simulation codebase, are openly available in the public repository at \url{https://github.com/Donzhumber/TPV.git}.
    \item \textbf{Author Contributions:} Prof. Humberto Bernal, Ph.D. is the sole author and contributed to all aspects of the conceptualization, mathematical proofs, econometric estimation, deep learning architecture, dynamic simulations, and drafting of the manuscript.
    \item \textbf{Declaration on the Use of Generative AI:} In the preparation of this manuscript, the author used generative AI tools (\emph{Cursor} and \emph{Antigravity}) for language copy-editing and assistance with software syntax in Lean~4, Stata, and Python. The author validated all results and takes full responsibility for the paper's contents.
    \item \textbf{Acknowledgments:} The author expresses sincere gratitude to academic colleagues and seminar participants at Universidad de los Andes and Universidad Colegio Mayor de Cundinamarca for invaluable feedback, methodological suggestions, and institutional support throughout the development of this research project.
\end{itemize}

\end{document}
